\documentclass{article}
\usepackage[paperwidth=180mm,paperheight=180mm,margin=16mm]{geometry}
\usepackage{amsmath,unicode-math,xcolor}
\setmathfont{XITS Math}
\usepackage{experimental/flat-core}
\pagestyle{empty}
\begin{document}
\noindent\textbf{LuaLaTeX: actual row boxes}\par\medskip
\noindent\textcolor{blue}{Blue: each separate row hbox.}\par
\noindent\textcolor{green!50!black}{Green: enclosing formula box.}
\par\bigskip
\noindent\textbf{1. Same example as the XeTeX preview}
\[
 \begingroup\fboxsep=3pt\fboxrule=.3pt
 \fcolorbox{green!50!black}{white}{$\displaystyle
 \FlatShowBoxes\begin{flatmatchedsplit}
  \left(a+b\\+c\right)^2_i
 \end{flatmatchedsplit}$}\endgroup
\]
\noindent Without outlines:
\[
 \begin{flatmatchedsplit}\left(a+b\\+c\right)^2_i\end{flatmatchedsplit}
\]
\noindent\textbf{2. A fraction and scripts on a spanning pair}
\[
 \begingroup\fboxsep=3pt\fboxrule=.3pt
 \fcolorbox{green!50!black}{white}{$\displaystyle
 \FlatShowBoxes\begin{flatmatchedsplit}
  f&=\left(a+\frac bc\\&+d\right)^2_i+z
 \end{flatmatchedsplit}$}\endgroup
\]
\noindent These are physically separate horizontal lists inside a vertical
box. Each row has its own height and depth. The prototype inserts at least
3pt between adjacent row boxes, with a minimum baseline distance of 15pt.
This spacing policy has not yet been matched to amsmath.
\end{document}
